\documentclass[12pt,a4paper]{article}
\usepackage[utf8]{inputenc}
\usepackage[T1]{fontenc}
\usepackage[brazil]{babel}
\usepackage[margin=2.5cm]{geometry}
\usepackage{xcolor}
\usepackage{booktabs}
\usepackage{longtable}
\usepackage{array}
\usepackage{parskip}
\usepackage{fancyhdr}
\usepackage{titlesec}
\usepackage{enumitem}

\definecolor{corOK}{HTML}{1a7a3c}
\definecolor{corERRO}{HTML}{c0392b}
\definecolor{corINFO}{HTML}{2471a3}
\definecolor{corVERIFICAR}{HTML}{d68910}
\definecolor{corFundoReprovado}{HTML}{fadbd8}
\definecolor{corFundoAprovado}{HTML}{d5f5e3}
\definecolor{cinzaClaro}{HTML}{f2f3f4}

\newcommand{\statusOK}[1]{\textcolor{corOK}{\textbf{[OK]}} #1}
\newcommand{\statusERRO}[1]{\textcolor{corERRO}{\textbf{[ERRO]}} #1}
\newcommand{\statusINFO}[1]{\textcolor{corINFO}{\textbf{[INFO]}} #1}
\newcommand{\statusVERIFICAR}[1]{\textcolor{corVERIFICAR}{\textbf{[VERIFICAR]}} #1}
\newcommand{\status}[1]{\texttt{#1}}

\pagestyle{fancy}
\fancyhf{}
\rhead{Verificador de Contestação -- Fase 1}
\rfoot{Página \thepage}
\lfoot{Gerado em 07/07/2026}

\titleformat{\section}{\large\bfseries}{}{0em}{}[\titlerule]
\titleformat{\subsection}{\normalsize\bfseries}{}{0em}{}

\setcounter{tocdepth}{1}
\setcounter{secnumdepth}{1}

\begin{document}

\begin{center}
  {\LARGE \textbf{Relatório Consolidado de Verificação de Contestações}}\\[0.4em]
  {\large Verificador de Contestação --- Fase 1}\\[0.2em]
  {\small Gerado em 07/07/2026 \quad|\quad 14 contestações verificadas}
\end{center}

\vspace{0.8em}

\noindent\colorbox{cinzaClaro}{%
  \begin{minipage}{\dimexpr\linewidth-2\fboxsep}
    \centering
    \vspace{0.4em}
    {\large \textbf{Resumo geral}}\\[0.3em]
    \textcolor{corOK}{\textbf{10 APROVADA(S)}} \quad|\quad \textcolor{corERRO}{\textbf{4 REPROVADA(S)}}
    \vspace{0.4em}
  \end{minipage}%
}

\vspace{1em}

\begin{longtable}{c p{6.5cm} c c c}
  \toprule
  \textbf{Nº} & \textbf{Autor} & \textbf{Veredito} & \textbf{Erros} & \textbf{Verif. manual} \\
  \midrule
  \endfirsthead
  \toprule
  \textbf{Nº} & \textbf{Autor} & \textbf{Veredito} & \textbf{Erros} & \textbf{Verif. manual} \\
  \midrule
  \endhead
  03 & ANTÔNIO PEREIRA LIMA & \textcolor{corOK}{\textbf{APROVADO}} & 0 & 0 \\
  04 & MARIA DAS GRAÇAS OLIVEIRA & \textcolor{corOK}{\textbf{APROVADO}} & 0 & 0 \\
  05 & SEBASTIÃO RODRIGUES DA COSTA & \textcolor{corERRO}{\textbf{REPROVADO}} & 1 & 0 \\
  06 & CARLOS ALBERTO FONSECA & \textcolor{corOK}{\textbf{APROVADO}} & 0 & 0 \\
  07 & JOSEFA MARIA DOS SANTOS & \textcolor{corOK}{\textbf{APROVADO}} & 0 & 0 \\
  08 & PAULO ROBERTO ALVES & \textcolor{corOK}{\textbf{APROVADO}} & 0 & 0 \\
  09 & ANA CLÁUDIA MENEZES & \textcolor{corOK}{\textbf{APROVADO}} & 0 & 0 \\
  10 & ROBERTO CARLOS DINIZ & \textcolor{corERRO}{\textbf{REPROVADO}} & 1 & 0 \\
  11 & FRANCISCO DE ASSIS MOURA & \textcolor{corOK}{\textbf{APROVADO}} & 0 & 0 \\
  12 & LUIZA HELENA BARROS & \textcolor{corERRO}{\textbf{REPROVADO}} & 1 & 0 \\
  13 & MARCOS VINÍCIUS TEIXEIRA & \textcolor{corOK}{\textbf{APROVADO}} & 0 & 0 \\
  14 & RAIMUNDO NONATO SILVA & \textcolor{corOK}{\textbf{APROVADO}} & 0 & 0 \\
  15 & TEREZA CRISTINA LOPES & \textcolor{corERRO}{\textbf{REPROVADO}} & 1 & 0 \\
  16 & EDSON LUÍS CAVALCANTE & \textcolor{corOK}{\textbf{APROVADO}} & 0 & 0 \\
  \bottomrule
\end{longtable}

\clearpage
\tableofcontents


\clearpage
\lhead{Contestação 03}
\section{Contestação 03 --- ANTÔNIO PEREIRA LIMA}
% -----------------------------------------------------------------------
% Cabeçalho
% -----------------------------------------------------------------------
\begin{center}
  {\small Processo n.º 5002233-11.2024.4.03.6301 \quad NB: 46/201.334.556-7 \quad DER: 18/05/2023}
\end{center}

\vspace{0.5em}

\noindent\colorbox{corFundoAprovado}{%
  \begin{minipage}{\dimexpr\linewidth-2\fboxsep}
    \centering
    \vspace{0.3em}
    {\large \textcolor{corOK}{\textbf{VEREDITO: APROVADO}}}\\
    Erros: \textbf{0} \quad|\quad Pontos de verificação manual: \textbf{0}
    \vspace{0.3em}
  \end{minipage}%
}

\vspace{1em}


% -----------------------------------------------------------------------
% Bloco 1 - Estrutural
% -----------------------------------------------------------------------
\subsection{Bloco 1 --- Estrutural}

\statusINFO{Tipo de endereçamento (JEF x Justiça Federal comum)}\\
Endereçamento identificado: \textbf{JEF (Juizado Especial Federal)}.\\
\textit{Evidência: ... SENHOR(A) DOUTOR(A) JUIZ(A) FEDERAL DA 8ª VARA FEDERAL DO JUIZADO ESPECIAL FEDERAL DA SUBSEÇÃO JUDICIÁRIA DE SÃO PAULO - SP Processo nº 500223...}

\medskip
\statusOK{Campo obrigatório: NB} --- NB presente.

\medskip
\statusOK{Campo obrigatório: DER} --- DER presente.

\medskip
\statusOK{Campo obrigatório: Número do processo} --- Número do processo presente.

\medskip
\statusOK{Campo obrigatório: Nome do autor} --- Nome do autor presente.\\
\textit{Evidência: ... a federal já qualificada nos autos do processo em epígrafe, que lhe move ANTÔNIO PEREIRA LIMA, vem, respeitosamente, à presença de Vossa Excelência, por...}

\medskip
\statusOK{Campo obrigatório: Tipo de benefício} --- Tipo de benefício presente.

\medskip
\statusOK{Marcações de edição não removidas (colchetes, placeholders, instruções, amarelo)}\\
Nenhuma marcação de edição remanescente detectada.

% -----------------------------------------------------------------------
% Bloco 2 - Preliminares
% -----------------------------------------------------------------------
\subsection{Bloco 2 --- Preliminares (presença)}

\begin{longtable}{p{8.5cm} p{4.5cm}}
  \toprule
  \textbf{Preliminar} & \textbf{Status} \\
  \midrule
  \endfirsthead
  \toprule
  \textbf{Preliminar} & \textbf{Status} \\
  \midrule
  \endhead
  Juízo 100\% Digital                           & \textcolor{corOK}{\textbf{PRESENTE}} \\
  Audiência de Conciliação                      & \textcolor{corOK}{\textbf{PRESENTE}} \\
  Renúncia aos 60 Salários-Mínimos              & \textcolor{corOK}{\textbf{PRESENTE}} \\
  Decadência                                    & AUSENTE \\
  Prescrição                                    & AUSENTE \\
  Coisa Julgada                                 & AUSENTE \\
  Litispendência                                & AUSENTE \\
  PPP não apresentado administrativamente       & AUSENTE \\
  Períodos reconhecidos administrativamente     & AUSENTE \\
  Petição inicial inepta                        & AUSENTE \\
  \bottomrule
\end{longtable}

\statusOK{Renúncia aos 60 SM obrigatória quando o juízo é o JEF}\\
Processo no JEF e preliminar de Renúncia aos 60 SM presente.

% -----------------------------------------------------------------------
% Bloco 6 - Redistribuição
% -----------------------------------------------------------------------
\subsection{Bloco 6 --- Redistribuição}

\statusOK{Indício de categoria vigilante/vigia/guarda/policial} --- Sem indício detectado.

\medskip
\statusOK{Indício de profissional/ambiente de saúde + agente biológico} --- Sem indício detectado.

\medskip
\statusOK{Indício de professor com período posterior a 1981} --- Sem indício detectado.

% -----------------------------------------------------------------------
% Teses
% -----------------------------------------------------------------------
\subsection{Teses presentes (Blocos 3/4/5 --- catalogação)}

\statusINFO{Catalogação das teses padronizadas presentes (CTN-/DIVESP- etc.)}\\
Nenhuma tese padronizada (CTN-/DIVESP-) encontrada no texto.

% -----------------------------------------------------------------------
% Verificações adiadas
% -----------------------------------------------------------------------
\subsection{Verificações adiadas para a fase do LLM}

\begin{itemize}
  \item \texttt{[3.x]} Compatibilidade tese × agente × período (marcos temporais)
  \item \texttt{[2.8]} Distinção entre as três situações de PPP não apresentado
  \item \texttt{[2.45]} Cabimento de decadência/prescrição com base em datas e tipo de ação
  \item \texttt{[1.4]} Cobertura dos períodos requeridos na petição inicial
\end{itemize}



\clearpage
\lhead{Contestação 04}
\section{Contestação 04 --- MARIA DAS GRAÇAS OLIVEIRA}
% -----------------------------------------------------------------------
% Cabeçalho
% -----------------------------------------------------------------------
\begin{center}
  {\small Processo n.º 5004410-22.2025.4.03.6100 \quad NB: 46/222.111.000-9 \quad DER: 12/04/2024}
\end{center}

\vspace{0.5em}

\noindent\colorbox{corFundoAprovado}{%
  \begin{minipage}{\dimexpr\linewidth-2\fboxsep}
    \centering
    \vspace{0.3em}
    {\large \textcolor{corOK}{\textbf{VEREDITO: APROVADO}}}\\
    Erros: \textbf{0} \quad|\quad Pontos de verificação manual: \textbf{0}
    \vspace{0.3em}
  \end{minipage}%
}

\vspace{1em}


% -----------------------------------------------------------------------
% Bloco 1 - Estrutural
% -----------------------------------------------------------------------
\subsection{Bloco 1 --- Estrutural}

\statusINFO{Tipo de endereçamento (JEF x Justiça Federal comum)}\\
Endereçamento identificado: \textbf{Justiça Federal comum (sem menção a JEF)}.\\
\textit{Evidência: ... EXCELENTÍSSIMO(A) SENHOR(A) DOUTOR(A) JUIZ(A) FEDERAL DA 3ª VARA FEDERAL DA SUBSEÇÃO JUDICIÁRIA DE SÃO PAULO - SP Processo nº 500441...}

\medskip
\statusOK{Campo obrigatório: NB} --- NB presente.

\medskip
\statusOK{Campo obrigatório: DER} --- DER presente.

\medskip
\statusOK{Campo obrigatório: Número do processo} --- Número do processo presente.

\medskip
\statusOK{Campo obrigatório: Nome do autor} --- Nome do autor presente.\\
\textit{Evidência: ... a federal já qualificada nos autos do processo em epígrafe, que lhe move MARIA DAS GRAÇAS OLIVEIRA, vem, respeitosamente, à presença de Vossa Excelência, por...}

\medskip
\statusOK{Campo obrigatório: Tipo de benefício} --- Tipo de benefício presente.

\medskip
\statusOK{Marcações de edição não removidas (colchetes, placeholders, instruções, amarelo)}\\
Nenhuma marcação de edição remanescente detectada.

% -----------------------------------------------------------------------
% Bloco 2 - Preliminares
% -----------------------------------------------------------------------
\subsection{Bloco 2 --- Preliminares (presença)}

\begin{longtable}{p{8.5cm} p{4.5cm}}
  \toprule
  \textbf{Preliminar} & \textbf{Status} \\
  \midrule
  \endfirsthead
  \toprule
  \textbf{Preliminar} & \textbf{Status} \\
  \midrule
  \endhead
  Juízo 100\% Digital                           & \textcolor{corOK}{\textbf{PRESENTE}} \\
  Audiência de Conciliação                      & \textcolor{corOK}{\textbf{PRESENTE}} \\
  Renúncia aos 60 Salários-Mínimos              & AUSENTE \\
  Decadência                                    & AUSENTE \\
  Prescrição                                    & AUSENTE \\
  Coisa Julgada                                 & AUSENTE \\
  Litispendência                                & AUSENTE \\
  PPP não apresentado administrativamente       & AUSENTE \\
  Períodos reconhecidos administrativamente     & AUSENTE \\
  Petição inicial inepta                        & AUSENTE \\
  \bottomrule
\end{longtable}

\statusINFO{Renúncia aos 60 SM obrigatória quando o juízo é o JEF}\\
Não é JEF; a exigência de renúncia aos 60 SM não se aplica.

% -----------------------------------------------------------------------
% Bloco 6 - Redistribuição
% -----------------------------------------------------------------------
\subsection{Bloco 6 --- Redistribuição}

\statusOK{Indício de categoria vigilante/vigia/guarda/policial} --- Sem indício detectado.

\medskip
\statusOK{Indício de profissional/ambiente de saúde + agente biológico} --- Sem indício detectado.

\medskip
\statusOK{Indício de professor com período posterior a 1981} --- Sem indício detectado.

% -----------------------------------------------------------------------
% Teses
% -----------------------------------------------------------------------
\subsection{Teses presentes (Blocos 3/4/5 --- catalogação)}

\statusINFO{Catalogação das teses padronizadas presentes (CTN-/DIVESP- etc.)}\\
Nenhuma tese padronizada (CTN-/DIVESP-) encontrada no texto.

% -----------------------------------------------------------------------
% Verificações adiadas
% -----------------------------------------------------------------------
\subsection{Verificações adiadas para a fase do LLM}

\begin{itemize}
  \item \texttt{[3.x]} Compatibilidade tese × agente × período (marcos temporais)
  \item \texttt{[2.8]} Distinção entre as três situações de PPP não apresentado
  \item \texttt{[2.45]} Cabimento de decadência/prescrição com base em datas e tipo de ação
  \item \texttt{[1.4]} Cobertura dos períodos requeridos na petição inicial
\end{itemize}



\clearpage
\lhead{Contestação 05}
\section{Contestação 05 --- SEBASTIÃO RODRIGUES DA COSTA}
% -----------------------------------------------------------------------
% Cabeçalho
% -----------------------------------------------------------------------
\begin{center}
  {\small Processo n.º 5006677-33.2024.4.03.6183 \quad NB: 46/333.222.111-4 \quad DER: 10/09/2023}
\end{center}

\vspace{0.5em}

\noindent\colorbox{corFundoReprovado}{%
  \begin{minipage}{\dimexpr\linewidth-2\fboxsep}
    \centering
    \vspace{0.3em}
    {\large \textcolor{corERRO}{\textbf{VEREDITO: REPROVADO}}}\\
    Erros: \textbf{1} \quad|\quad Pontos de verificação manual: \textbf{0}
    \vspace{0.3em}
  \end{minipage}%
}

\vspace{1em}


% -----------------------------------------------------------------------
% Bloco 1 - Estrutural
% -----------------------------------------------------------------------
\subsection{Bloco 1 --- Estrutural}

\statusINFO{Tipo de endereçamento (JEF x Justiça Federal comum)}\\
Endereçamento identificado: \textbf{JEF (Juizado Especial Federal)}.\\
\textit{Evidência: ... SENHOR(A) DOUTOR(A) JUIZ(A) FEDERAL DA 10ª VARA FEDERAL DO JUIZADO ESPECIAL FEDERAL DA SUBSEÇÃO JUDICIÁRIA DE SÃO PAULO - SP Processo nº 500667...}

\medskip
\statusOK{Campo obrigatório: NB} --- NB presente.

\medskip
\statusOK{Campo obrigatório: DER} --- DER presente.

\medskip
\statusOK{Campo obrigatório: Número do processo} --- Número do processo presente.

\medskip
\statusOK{Campo obrigatório: Nome do autor} --- Nome do autor presente.\\
\textit{Evidência: ... a federal já qualificada nos autos do processo em epígrafe, que lhe move SEBASTIÃO RODRIGUES DA COSTA, vem, respeitosamente, à presença de Vossa Excelência, por...}

\medskip
\statusOK{Campo obrigatório: Tipo de benefício} --- Tipo de benefício presente.

\medskip
\statusERRO{Marcações de edição não removidas (colchetes, placeholders, instruções, amarelo)}\\
Há 2 marcação(ões) de edição não removida(s) na minuta --- devem ser eliminadas antes do envio.\\
\textit{Evidência: instrução remanescente: "...r o ato de concessão de seu benefício. [VERIFICAR: deixar a preliminar de decadência some..." | instrução remanescente: "...concessão de seu benefício. [VERIFICAR: deixar a preliminar de decadência somente quando contabiliz..."}

% -----------------------------------------------------------------------
% Bloco 2 - Preliminares
% -----------------------------------------------------------------------
\subsection{Bloco 2 --- Preliminares (presença)}

\begin{longtable}{p{8.5cm} p{4.5cm}}
  \toprule
  \textbf{Preliminar} & \textbf{Status} \\
  \midrule
  \endfirsthead
  \toprule
  \textbf{Preliminar} & \textbf{Status} \\
  \midrule
  \endhead
  Juízo 100\% Digital                           & \textcolor{corOK}{\textbf{PRESENTE}} \\
  Audiência de Conciliação                      & \textcolor{corOK}{\textbf{PRESENTE}} \\
  Renúncia aos 60 Salários-Mínimos              & \textcolor{corOK}{\textbf{PRESENTE}} \\
  Decadência                                    & \textcolor{corOK}{\textbf{PRESENTE}} \\
  Prescrição                                    & AUSENTE \\
  Coisa Julgada                                 & AUSENTE \\
  Litispendência                                & AUSENTE \\
  PPP não apresentado administrativamente       & AUSENTE \\
  Períodos reconhecidos administrativamente     & AUSENTE \\
  Petição inicial inepta                        & AUSENTE \\
  \bottomrule
\end{longtable}

\statusOK{Renúncia aos 60 SM obrigatória quando o juízo é o JEF}\\
Processo no JEF e preliminar de Renúncia aos 60 SM presente.

% -----------------------------------------------------------------------
% Bloco 6 - Redistribuição
% -----------------------------------------------------------------------
\subsection{Bloco 6 --- Redistribuição}

\statusOK{Indício de categoria vigilante/vigia/guarda/policial} --- Sem indício detectado.

\medskip
\statusOK{Indício de profissional/ambiente de saúde + agente biológico} --- Sem indício detectado.

\medskip
\statusOK{Indício de professor com período posterior a 1981} --- Sem indício detectado.

% -----------------------------------------------------------------------
% Teses
% -----------------------------------------------------------------------
\subsection{Teses presentes (Blocos 3/4/5 --- catalogação)}

\statusINFO{Catalogação das teses padronizadas presentes (CTN-/DIVESP- etc.)}\\
Nenhuma tese padronizada (CTN-/DIVESP-) encontrada no texto.

% -----------------------------------------------------------------------
% Verificações adiadas
% -----------------------------------------------------------------------
\subsection{Verificações adiadas para a fase do LLM}

\begin{itemize}
  \item \texttt{[3.x]} Compatibilidade tese × agente × período (marcos temporais)
  \item \texttt{[2.8]} Distinção entre as três situações de PPP não apresentado
  \item \texttt{[2.45]} Cabimento de decadência/prescrição com base em datas e tipo de ação
  \item \texttt{[1.4]} Cobertura dos períodos requeridos na petição inicial
\end{itemize}



\clearpage
\lhead{Contestação 06}
\section{Contestação 06 --- CARLOS ALBERTO FONSECA}
% -----------------------------------------------------------------------
% Cabeçalho
% -----------------------------------------------------------------------
\begin{center}
  {\small Processo n.º 5007788-44.2025.4.03.6183 \quad NB: 46/444.333.222-1 \quad DER: 22/01/2025}
\end{center}

\vspace{0.5em}

\noindent\colorbox{corFundoAprovado}{%
  \begin{minipage}{\dimexpr\linewidth-2\fboxsep}
    \centering
    \vspace{0.3em}
    {\large \textcolor{corOK}{\textbf{VEREDITO: APROVADO}}}\\
    Erros: \textbf{0} \quad|\quad Pontos de verificação manual: \textbf{0}
    \vspace{0.3em}
  \end{minipage}%
}

\vspace{1em}


% -----------------------------------------------------------------------
% Bloco 1 - Estrutural
% -----------------------------------------------------------------------
\subsection{Bloco 1 --- Estrutural}

\statusINFO{Tipo de endereçamento (JEF x Justiça Federal comum)}\\
Endereçamento identificado: \textbf{JEF (Juizado Especial Federal)}.\\
\textit{Evidência: ... SENHOR(A) DOUTOR(A) JUIZ(A) FEDERAL DA 6ª VARA FEDERAL DO JUIZADO ESPECIAL FEDERAL DA SUBSEÇÃO JUDICIÁRIA DE SÃO PAULO - SP Processo nº 500778...}

\medskip
\statusOK{Campo obrigatório: NB} --- NB presente.

\medskip
\statusOK{Campo obrigatório: DER} --- DER presente.

\medskip
\statusOK{Campo obrigatório: Número do processo} --- Número do processo presente.

\medskip
\statusOK{Campo obrigatório: Nome do autor} --- Nome do autor presente.\\
\textit{Evidência: ... a federal já qualificada nos autos do processo em epígrafe, que lhe move CARLOS ALBERTO FONSECA, vem, respeitosamente, à presença de Vossa Excelência, por...}

\medskip
\statusOK{Campo obrigatório: Tipo de benefício} --- Tipo de benefício presente.

\medskip
\statusOK{Marcações de edição não removidas (colchetes, placeholders, instruções, amarelo)}\\
Nenhuma marcação de edição remanescente detectada.

% -----------------------------------------------------------------------
% Bloco 2 - Preliminares
% -----------------------------------------------------------------------
\subsection{Bloco 2 --- Preliminares (presença)}

\begin{longtable}{p{8.5cm} p{4.5cm}}
  \toprule
  \textbf{Preliminar} & \textbf{Status} \\
  \midrule
  \endfirsthead
  \toprule
  \textbf{Preliminar} & \textbf{Status} \\
  \midrule
  \endhead
  Juízo 100\% Digital                           & \textcolor{corOK}{\textbf{PRESENTE}} \\
  Audiência de Conciliação                      & \textcolor{corOK}{\textbf{PRESENTE}} \\
  Renúncia aos 60 Salários-Mínimos              & \textcolor{corOK}{\textbf{PRESENTE}} \\
  Decadência                                    & AUSENTE \\
  Prescrição                                    & AUSENTE \\
  Coisa Julgada                                 & AUSENTE \\
  Litispendência                                & AUSENTE \\
  PPP não apresentado administrativamente       & AUSENTE \\
  Períodos reconhecidos administrativamente     & AUSENTE \\
  Petição inicial inepta                        & AUSENTE \\
  \bottomrule
\end{longtable}

\statusOK{Renúncia aos 60 SM obrigatória quando o juízo é o JEF}\\
Processo no JEF e preliminar de Renúncia aos 60 SM presente.

% -----------------------------------------------------------------------
% Bloco 6 - Redistribuição
% -----------------------------------------------------------------------
\subsection{Bloco 6 --- Redistribuição}

\statusOK{Indício de categoria vigilante/vigia/guarda/policial} --- Sem indício detectado.

\medskip
\statusOK{Indício de profissional/ambiente de saúde + agente biológico} --- Sem indício detectado.

\medskip
\statusOK{Indício de professor com período posterior a 1981} --- Sem indício detectado.

% -----------------------------------------------------------------------
% Teses
% -----------------------------------------------------------------------
\subsection{Teses presentes (Blocos 3/4/5 --- catalogação)}

\statusINFO{Catalogação das teses padronizadas presentes (CTN-/DIVESP- etc.)}\\
Nenhuma tese padronizada (CTN-/DIVESP-) encontrada no texto.

% -----------------------------------------------------------------------
% Verificações adiadas
% -----------------------------------------------------------------------
\subsection{Verificações adiadas para a fase do LLM}

\begin{itemize}
  \item \texttt{[3.x]} Compatibilidade tese × agente × período (marcos temporais)
  \item \texttt{[2.8]} Distinção entre as três situações de PPP não apresentado
  \item \texttt{[2.45]} Cabimento de decadência/prescrição com base em datas e tipo de ação
  \item \texttt{[1.4]} Cobertura dos períodos requeridos na petição inicial
\end{itemize}



\clearpage
\lhead{Contestação 07}
\section{Contestação 07 --- JOSEFA MARIA DOS SANTOS}
% -----------------------------------------------------------------------
% Cabeçalho
% -----------------------------------------------------------------------
\begin{center}
  {\small Processo n.º 5008899-55.2025.4.03.6301 \quad NB: 46/555.444.333-8 \quad DER: 30/06/2024}
\end{center}

\vspace{0.5em}

\noindent\colorbox{corFundoAprovado}{%
  \begin{minipage}{\dimexpr\linewidth-2\fboxsep}
    \centering
    \vspace{0.3em}
    {\large \textcolor{corOK}{\textbf{VEREDITO: APROVADO}}}\\
    Erros: \textbf{0} \quad|\quad Pontos de verificação manual: \textbf{0}
    \vspace{0.3em}
  \end{minipage}%
}

\vspace{1em}


% -----------------------------------------------------------------------
% Bloco 1 - Estrutural
% -----------------------------------------------------------------------
\subsection{Bloco 1 --- Estrutural}

\statusINFO{Tipo de endereçamento (JEF x Justiça Federal comum)}\\
Endereçamento identificado: \textbf{JEF (Juizado Especial Federal)}.\\
\textit{Evidência: ... SENHOR(A) DOUTOR(A) JUIZ(A) FEDERAL DA 4ª VARA FEDERAL DO JUIZADO ESPECIAL FEDERAL DA SUBSEÇÃO JUDICIÁRIA DE SÃO PAULO - SP Processo nº 500889...}

\medskip
\statusOK{Campo obrigatório: NB} --- NB presente.

\medskip
\statusOK{Campo obrigatório: DER} --- DER presente.

\medskip
\statusOK{Campo obrigatório: Número do processo} --- Número do processo presente.

\medskip
\statusOK{Campo obrigatório: Nome do autor} --- Nome do autor presente.\\
\textit{Evidência: ... a federal já qualificada nos autos do processo em epígrafe, que lhe move JOSEFA MARIA DOS SANTOS, vem, respeitosamente, à presença de Vossa Excelência, por...}

\medskip
\statusOK{Campo obrigatório: Tipo de benefício} --- Tipo de benefício presente.

\medskip
\statusOK{Marcações de edição não removidas (colchetes, placeholders, instruções, amarelo)}\\
Nenhuma marcação de edição remanescente detectada.

% -----------------------------------------------------------------------
% Bloco 2 - Preliminares
% -----------------------------------------------------------------------
\subsection{Bloco 2 --- Preliminares (presença)}

\begin{longtable}{p{8.5cm} p{4.5cm}}
  \toprule
  \textbf{Preliminar} & \textbf{Status} \\
  \midrule
  \endfirsthead
  \toprule
  \textbf{Preliminar} & \textbf{Status} \\
  \midrule
  \endhead
  Juízo 100\% Digital                           & \textcolor{corOK}{\textbf{PRESENTE}} \\
  Audiência de Conciliação                      & \textcolor{corOK}{\textbf{PRESENTE}} \\
  Renúncia aos 60 Salários-Mínimos              & \textcolor{corOK}{\textbf{PRESENTE}} \\
  Decadência                                    & AUSENTE \\
  Prescrição                                    & AUSENTE \\
  Coisa Julgada                                 & AUSENTE \\
  Litispendência                                & AUSENTE \\
  PPP não apresentado administrativamente       & AUSENTE \\
  Períodos reconhecidos administrativamente     & AUSENTE \\
  Petição inicial inepta                        & AUSENTE \\
  \bottomrule
\end{longtable}

\statusOK{Renúncia aos 60 SM obrigatória quando o juízo é o JEF}\\
Processo no JEF e preliminar de Renúncia aos 60 SM presente.

% -----------------------------------------------------------------------
% Bloco 6 - Redistribuição
% -----------------------------------------------------------------------
\subsection{Bloco 6 --- Redistribuição}

\statusOK{Indício de categoria vigilante/vigia/guarda/policial} --- Sem indício detectado.

\medskip
\statusOK{Indício de profissional/ambiente de saúde + agente biológico} --- Sem indício detectado.

\medskip
\statusOK{Indício de professor com período posterior a 1981} --- Sem indício detectado.

% -----------------------------------------------------------------------
% Teses
% -----------------------------------------------------------------------
\subsection{Teses presentes (Blocos 3/4/5 --- catalogação)}

\statusINFO{Catalogação das teses padronizadas presentes (CTN-/DIVESP- etc.)}\\
Nenhuma tese padronizada (CTN-/DIVESP-) encontrada no texto.

% -----------------------------------------------------------------------
% Verificações adiadas
% -----------------------------------------------------------------------
\subsection{Verificações adiadas para a fase do LLM}

\begin{itemize}
  \item \texttt{[3.x]} Compatibilidade tese × agente × período (marcos temporais)
  \item \texttt{[2.8]} Distinção entre as três situações de PPP não apresentado
  \item \texttt{[2.45]} Cabimento de decadência/prescrição com base em datas e tipo de ação
  \item \texttt{[1.4]} Cobertura dos períodos requeridos na petição inicial
\end{itemize}



\clearpage
\lhead{Contestação 08}
\section{Contestação 08 --- PAULO ROBERTO ALVES}
% -----------------------------------------------------------------------
% Cabeçalho
% -----------------------------------------------------------------------
\begin{center}
  {\small Processo n.º 5009900-66.2025.4.03.6183 \quad NB: 46/666.555.444-5 \quad DER: 14/03/2024}
\end{center}

\vspace{0.5em}

\noindent\colorbox{corFundoAprovado}{%
  \begin{minipage}{\dimexpr\linewidth-2\fboxsep}
    \centering
    \vspace{0.3em}
    {\large \textcolor{corOK}{\textbf{VEREDITO: APROVADO}}}\\
    Erros: \textbf{0} \quad|\quad Pontos de verificação manual: \textbf{0}
    \vspace{0.3em}
  \end{minipage}%
}

\vspace{1em}


% -----------------------------------------------------------------------
% Bloco 1 - Estrutural
% -----------------------------------------------------------------------
\subsection{Bloco 1 --- Estrutural}

\statusINFO{Tipo de endereçamento (JEF x Justiça Federal comum)}\\
Endereçamento identificado: \textbf{JEF (Juizado Especial Federal)}.\\
\textit{Evidência: ... SENHOR(A) DOUTOR(A) JUIZ(A) FEDERAL DA 7ª VARA FEDERAL DO JUIZADO ESPECIAL FEDERAL DA SUBSEÇÃO JUDICIÁRIA DE SÃO PAULO - SP Processo nº 500990...}

\medskip
\statusOK{Campo obrigatório: NB} --- NB presente.

\medskip
\statusOK{Campo obrigatório: DER} --- DER presente.

\medskip
\statusOK{Campo obrigatório: Número do processo} --- Número do processo presente.

\medskip
\statusOK{Campo obrigatório: Nome do autor} --- Nome do autor presente.\\
\textit{Evidência: ... a federal já qualificada nos autos do processo em epígrafe, que lhe move PAULO ROBERTO ALVES, vem, respeitosamente, à presença de Vossa Excelência, por...}

\medskip
\statusOK{Campo obrigatório: Tipo de benefício} --- Tipo de benefício presente.

\medskip
\statusOK{Marcações de edição não removidas (colchetes, placeholders, instruções, amarelo)}\\
Nenhuma marcação de edição remanescente detectada.

% -----------------------------------------------------------------------
% Bloco 2 - Preliminares
% -----------------------------------------------------------------------
\subsection{Bloco 2 --- Preliminares (presença)}

\begin{longtable}{p{8.5cm} p{4.5cm}}
  \toprule
  \textbf{Preliminar} & \textbf{Status} \\
  \midrule
  \endfirsthead
  \toprule
  \textbf{Preliminar} & \textbf{Status} \\
  \midrule
  \endhead
  Juízo 100\% Digital                           & \textcolor{corOK}{\textbf{PRESENTE}} \\
  Audiência de Conciliação                      & \textcolor{corOK}{\textbf{PRESENTE}} \\
  Renúncia aos 60 Salários-Mínimos              & \textcolor{corOK}{\textbf{PRESENTE}} \\
  Decadência                                    & AUSENTE \\
  Prescrição                                    & AUSENTE \\
  Coisa Julgada                                 & AUSENTE \\
  Litispendência                                & AUSENTE \\
  PPP não apresentado administrativamente       & AUSENTE \\
  Períodos reconhecidos administrativamente     & AUSENTE \\
  Petição inicial inepta                        & AUSENTE \\
  \bottomrule
\end{longtable}

\statusOK{Renúncia aos 60 SM obrigatória quando o juízo é o JEF}\\
Processo no JEF e preliminar de Renúncia aos 60 SM presente.

% -----------------------------------------------------------------------
% Bloco 6 - Redistribuição
% -----------------------------------------------------------------------
\subsection{Bloco 6 --- Redistribuição}

\statusINFO{Indício de categoria vigilante/vigia/guarda/policial}\\
\textcolor{corVERIFICAR}{\textbf{Redistribuição sugerida --- flag VIGILANTE}} (termo(s): vigilante).
\\
\textit{Evidência: aposentadoria especial, alegando exercício da atividade de VIGILANTE, com uso de arma de fogo, no período de 02/05/2008 a 18/07/...}

\medskip
\statusOK{Indício de profissional/ambiente de saúde + agente biológico} --- Sem indício detectado.

\medskip
\statusOK{Indício de professor com período posterior a 1981} --- Sem indício detectado.

% -----------------------------------------------------------------------
% Teses
% -----------------------------------------------------------------------
\subsection{Teses presentes (Blocos 3/4/5 --- catalogação)}

\statusINFO{Catalogação das teses padronizadas presentes (CTN-/DIVESP- etc.)}\\
Nenhuma tese padronizada (CTN-/DIVESP-) encontrada no texto.

% -----------------------------------------------------------------------
% Verificações adiadas
% -----------------------------------------------------------------------
\subsection{Verificações adiadas para a fase do LLM}

\begin{itemize}
  \item \texttt{[3.x]} Compatibilidade tese × agente × período (marcos temporais)
  \item \texttt{[2.8]} Distinção entre as três situações de PPP não apresentado
  \item \texttt{[2.45]} Cabimento de decadência/prescrição com base em datas e tipo de ação
  \item \texttt{[1.4]} Cobertura dos períodos requeridos na petição inicial
\end{itemize}



\clearpage
\lhead{Contestação 09}
\section{Contestação 09 --- ANA CLÁUDIA MENEZES}
% -----------------------------------------------------------------------
% Cabeçalho
% -----------------------------------------------------------------------
\begin{center}
  {\small Processo n.º 5010011-77.2025.4.03.6301 \quad NB: 46/777.666.555-2 \quad DER: 05/02/2024}
\end{center}

\vspace{0.5em}

\noindent\colorbox{corFundoAprovado}{%
  \begin{minipage}{\dimexpr\linewidth-2\fboxsep}
    \centering
    \vspace{0.3em}
    {\large \textcolor{corOK}{\textbf{VEREDITO: APROVADO}}}\\
    Erros: \textbf{0} \quad|\quad Pontos de verificação manual: \textbf{0}
    \vspace{0.3em}
  \end{minipage}%
}

\vspace{1em}


% -----------------------------------------------------------------------
% Bloco 1 - Estrutural
% -----------------------------------------------------------------------
\subsection{Bloco 1 --- Estrutural}

\statusINFO{Tipo de endereçamento (JEF x Justiça Federal comum)}\\
Endereçamento identificado: \textbf{JEF (Juizado Especial Federal)}.\\
\textit{Evidência: ... SENHOR(A) DOUTOR(A) JUIZ(A) FEDERAL DA 9ª VARA FEDERAL DO JUIZADO ESPECIAL FEDERAL DA SUBSEÇÃO JUDICIÁRIA DE SÃO PAULO - SP Processo nº 501001...}

\medskip
\statusOK{Campo obrigatório: NB} --- NB presente.

\medskip
\statusOK{Campo obrigatório: DER} --- DER presente.

\medskip
\statusOK{Campo obrigatório: Número do processo} --- Número do processo presente.

\medskip
\statusOK{Campo obrigatório: Nome do autor} --- Nome do autor presente.\\
\textit{Evidência: ... a federal já qualificada nos autos do processo em epígrafe, que lhe move ANA CLÁUDIA MENEZES, vem, respeitosamente, à presença de Vossa Excelência, por...}

\medskip
\statusOK{Campo obrigatório: Tipo de benefício} --- Tipo de benefício presente.

\medskip
\statusOK{Marcações de edição não removidas (colchetes, placeholders, instruções, amarelo)}\\
Nenhuma marcação de edição remanescente detectada.

% -----------------------------------------------------------------------
% Bloco 2 - Preliminares
% -----------------------------------------------------------------------
\subsection{Bloco 2 --- Preliminares (presença)}

\begin{longtable}{p{8.5cm} p{4.5cm}}
  \toprule
  \textbf{Preliminar} & \textbf{Status} \\
  \midrule
  \endfirsthead
  \toprule
  \textbf{Preliminar} & \textbf{Status} \\
  \midrule
  \endhead
  Juízo 100\% Digital                           & \textcolor{corOK}{\textbf{PRESENTE}} \\
  Audiência de Conciliação                      & \textcolor{corOK}{\textbf{PRESENTE}} \\
  Renúncia aos 60 Salários-Mínimos              & \textcolor{corOK}{\textbf{PRESENTE}} \\
  Decadência                                    & AUSENTE \\
  Prescrição                                    & AUSENTE \\
  Coisa Julgada                                 & AUSENTE \\
  Litispendência                                & AUSENTE \\
  PPP não apresentado administrativamente       & AUSENTE \\
  Períodos reconhecidos administrativamente     & AUSENTE \\
  Petição inicial inepta                        & AUSENTE \\
  \bottomrule
\end{longtable}

\statusOK{Renúncia aos 60 SM obrigatória quando o juízo é o JEF}\\
Processo no JEF e preliminar de Renúncia aos 60 SM presente.

% -----------------------------------------------------------------------
% Bloco 6 - Redistribuição
% -----------------------------------------------------------------------
\subsection{Bloco 6 --- Redistribuição}

\statusOK{Indício de categoria vigilante/vigia/guarda/policial} --- Sem indício detectado.

\medskip
\statusINFO{Indício de profissional/ambiente de saúde + agente biológico}\\
\textcolor{corVERIFICAR}{\textbf{Redistribuição sugerida --- flag SAÚDE}} (indício(s): auxiliar de enfermagem, hospital, ambiente hospitalar + agente biológico).
\\
\textit{Evidência: cessão do benefício de aposentadoria especial, na função de AUXILIAR DE ENFERMAGEM em ambiente hospitalar, alegando exposição a agente biológi...}

\medskip
\statusOK{Indício de professor com período posterior a 1981} --- Sem indício detectado.

% -----------------------------------------------------------------------
% Teses
% -----------------------------------------------------------------------
\subsection{Teses presentes (Blocos 3/4/5 --- catalogação)}

\statusINFO{Catalogação das teses padronizadas presentes (CTN-/DIVESP- etc.)}\\
Nenhuma tese padronizada (CTN-/DIVESP-) encontrada no texto.

% -----------------------------------------------------------------------
% Verificações adiadas
% -----------------------------------------------------------------------
\subsection{Verificações adiadas para a fase do LLM}

\begin{itemize}
  \item \texttt{[3.x]} Compatibilidade tese × agente × período (marcos temporais)
  \item \texttt{[2.8]} Distinção entre as três situações de PPP não apresentado
  \item \texttt{[2.45]} Cabimento de decadência/prescrição com base em datas e tipo de ação
  \item \texttt{[1.4]} Cobertura dos períodos requeridos na petição inicial
\end{itemize}



\clearpage
\lhead{Contestação 10}
\section{Contestação 10 --- ROBERTO CARLOS DINIZ}
% -----------------------------------------------------------------------
% Cabeçalho
% -----------------------------------------------------------------------
\begin{center}
  {\small Processo n.º 5011122-88.2025.4.03.6183 \quad NB: 46/888.777.666-3 \quad DER: 20/11/2023}
\end{center}

\vspace{0.5em}

\noindent\colorbox{corFundoReprovado}{%
  \begin{minipage}{\dimexpr\linewidth-2\fboxsep}
    \centering
    \vspace{0.3em}
    {\large \textcolor{corERRO}{\textbf{VEREDITO: REPROVADO}}}\\
    Erros: \textbf{1} \quad|\quad Pontos de verificação manual: \textbf{0}
    \vspace{0.3em}
  \end{minipage}%
}

\vspace{1em}


% -----------------------------------------------------------------------
% Bloco 1 - Estrutural
% -----------------------------------------------------------------------
\subsection{Bloco 1 --- Estrutural}

\statusINFO{Tipo de endereçamento (JEF x Justiça Federal comum)}\\
Endereçamento identificado: \textbf{JEF (Juizado Especial Federal)}.\\
\textit{Evidência: ... SENHOR(A) DOUTOR(A) JUIZ(A) FEDERAL DA 2ª VARA FEDERAL DO JUIZADO ESPECIAL FEDERAL DA SUBSEÇÃO JUDICIÁRIA DE SÃO PAULO - SP Processo nº 501112...}

\medskip
\statusOK{Campo obrigatório: NB} --- NB presente.

\medskip
\statusOK{Campo obrigatório: DER} --- DER presente.

\medskip
\statusOK{Campo obrigatório: Número do processo} --- Número do processo presente.

\medskip
\statusOK{Campo obrigatório: Nome do autor} --- Nome do autor presente.\\
\textit{Evidência: ... a federal já qualificada nos autos do processo em epígrafe, que lhe move ROBERTO CARLOS DINIZ, vem, respeitosamente, à presença de Vossa Excelência, por...}

\medskip
\statusERRO{Campo obrigatório: Tipo de benefício}\\
\textbf{Tipo de benefício ausente na minuta.}

\medskip
\statusOK{Marcações de edição não removidas (colchetes, placeholders, instruções, amarelo)}\\
Nenhuma marcação de edição remanescente detectada.

% -----------------------------------------------------------------------
% Bloco 2 - Preliminares
% -----------------------------------------------------------------------
\subsection{Bloco 2 --- Preliminares (presença)}

\begin{longtable}{p{8.5cm} p{4.5cm}}
  \toprule
  \textbf{Preliminar} & \textbf{Status} \\
  \midrule
  \endfirsthead
  \toprule
  \textbf{Preliminar} & \textbf{Status} \\
  \midrule
  \endhead
  Juízo 100\% Digital                           & \textcolor{corOK}{\textbf{PRESENTE}} \\
  Audiência de Conciliação                      & \textcolor{corOK}{\textbf{PRESENTE}} \\
  Renúncia aos 60 Salários-Mínimos              & \textcolor{corOK}{\textbf{PRESENTE}} \\
  Decadência                                    & AUSENTE \\
  Prescrição                                    & AUSENTE \\
  Coisa Julgada                                 & AUSENTE \\
  Litispendência                                & AUSENTE \\
  PPP não apresentado administrativamente       & AUSENTE \\
  Períodos reconhecidos administrativamente     & AUSENTE \\
  Petição inicial inepta                        & AUSENTE \\
  \bottomrule
\end{longtable}

\statusOK{Renúncia aos 60 SM obrigatória quando o juízo é o JEF}\\
Processo no JEF e preliminar de Renúncia aos 60 SM presente.

% -----------------------------------------------------------------------
% Bloco 6 - Redistribuição
% -----------------------------------------------------------------------
\subsection{Bloco 6 --- Redistribuição}

\statusOK{Indício de categoria vigilante/vigia/guarda/policial} --- Sem indício detectado.

\medskip
\statusOK{Indício de profissional/ambiente de saúde + agente biológico} --- Sem indício detectado.

\medskip
\statusINFO{Indício de professor com período posterior a 1981}\\
\textcolor{corVERIFICAR}{\textbf{Redistribuição sugerida --- flag PROFESSOR}} (professor com indício de período posterior a 1981).
\\
\textit{Evidência: arte autora o reconhecimento de tempo especial na função de PROFESSOR, em exercício de magistério, no período de 01/02/1985 a 31/...}

% -----------------------------------------------------------------------
% Teses
% -----------------------------------------------------------------------
\subsection{Teses presentes (Blocos 3/4/5 --- catalogação)}

\statusINFO{Catalogação das teses padronizadas presentes (CTN-/DIVESP- etc.)}\\
Nenhuma tese padronizada (CTN-/DIVESP-) encontrada no texto.

% -----------------------------------------------------------------------
% Verificações adiadas
% -----------------------------------------------------------------------
\subsection{Verificações adiadas para a fase do LLM}

\begin{itemize}
  \item \texttt{[3.x]} Compatibilidade tese × agente × período (marcos temporais)
  \item \texttt{[2.8]} Distinção entre as três situações de PPP não apresentado
  \item \texttt{[2.45]} Cabimento de decadência/prescrição com base em datas e tipo de ação
  \item \texttt{[1.4]} Cobertura dos períodos requeridos na petição inicial
\end{itemize}



\clearpage
\lhead{Contestação 11}
\section{Contestação 11 --- FRANCISCO DE ASSIS MOURA}
% -----------------------------------------------------------------------
% Cabeçalho
% -----------------------------------------------------------------------
\begin{center}
  {\small Processo n.º 5012233-99.2025.4.03.6301 \quad NB: 46/999.888.777-0 \quad DER: 08/01/2024}
\end{center}

\vspace{0.5em}

\noindent\colorbox{corFundoAprovado}{%
  \begin{minipage}{\dimexpr\linewidth-2\fboxsep}
    \centering
    \vspace{0.3em}
    {\large \textcolor{corOK}{\textbf{VEREDITO: APROVADO}}}\\
    Erros: \textbf{0} \quad|\quad Pontos de verificação manual: \textbf{0}
    \vspace{0.3em}
  \end{minipage}%
}

\vspace{1em}


% -----------------------------------------------------------------------
% Bloco 1 - Estrutural
% -----------------------------------------------------------------------
\subsection{Bloco 1 --- Estrutural}

\statusINFO{Tipo de endereçamento (JEF x Justiça Federal comum)}\\
Endereçamento identificado: \textbf{JEF (Juizado Especial Federal)}.\\
\textit{Evidência: ... SENHOR(A) DOUTOR(A) JUIZ(A) FEDERAL DA 5ª VARA FEDERAL DO JUIZADO ESPECIAL FEDERAL DA SUBSEÇÃO JUDICIÁRIA DE SÃO PAULO - SP Processo nº 501223...}

\medskip
\statusOK{Campo obrigatório: NB} --- NB presente.

\medskip
\statusOK{Campo obrigatório: DER} --- DER presente.

\medskip
\statusOK{Campo obrigatório: Número do processo} --- Número do processo presente.

\medskip
\statusOK{Campo obrigatório: Nome do autor} --- Nome do autor presente.\\
\textit{Evidência: ... a federal já qualificada nos autos do processo em epígrafe, que lhe move FRANCISCO DE ASSIS MOURA, vem, respeitosamente, à presença de Vossa Excelência, por...}

\medskip
\statusOK{Campo obrigatório: Tipo de benefício} --- Tipo de benefício presente.

\medskip
\statusOK{Marcações de edição não removidas (colchetes, placeholders, instruções, amarelo)}\\
Nenhuma marcação de edição remanescente detectada.

% -----------------------------------------------------------------------
% Bloco 2 - Preliminares
% -----------------------------------------------------------------------
\subsection{Bloco 2 --- Preliminares (presença)}

\begin{longtable}{p{8.5cm} p{4.5cm}}
  \toprule
  \textbf{Preliminar} & \textbf{Status} \\
  \midrule
  \endfirsthead
  \toprule
  \textbf{Preliminar} & \textbf{Status} \\
  \midrule
  \endhead
  Juízo 100\% Digital                           & \textcolor{corOK}{\textbf{PRESENTE}} \\
  Audiência de Conciliação                      & \textcolor{corOK}{\textbf{PRESENTE}} \\
  Renúncia aos 60 Salários-Mínimos              & \textcolor{corOK}{\textbf{PRESENTE}} \\
  Decadência                                    & AUSENTE \\
  Prescrição                                    & AUSENTE \\
  Coisa Julgada                                 & AUSENTE \\
  Litispendência                                & AUSENTE \\
  PPP não apresentado administrativamente       & AUSENTE \\
  Períodos reconhecidos administrativamente     & AUSENTE \\
  Petição inicial inepta                        & AUSENTE \\
  \bottomrule
\end{longtable}

\statusOK{Renúncia aos 60 SM obrigatória quando o juízo é o JEF}\\
Processo no JEF e preliminar de Renúncia aos 60 SM presente.

% -----------------------------------------------------------------------
% Bloco 6 - Redistribuição
% -----------------------------------------------------------------------
\subsection{Bloco 6 --- Redistribuição}

\statusOK{Indício de categoria vigilante/vigia/guarda/policial} --- Sem indício detectado.

\medskip
\statusOK{Indício de profissional/ambiente de saúde + agente biológico} --- Sem indício detectado.

\medskip
\statusOK{Indício de professor com período posterior a 1981} --- Sem indício detectado.

% -----------------------------------------------------------------------
% Teses
% -----------------------------------------------------------------------
\subsection{Teses presentes (Blocos 3/4/5 --- catalogação)}

\statusINFO{Catalogação das teses padronizadas presentes (CTN-/DIVESP- etc.)}\\
Nenhuma tese padronizada (CTN-/DIVESP-) encontrada no texto.

% -----------------------------------------------------------------------
% Verificações adiadas
% -----------------------------------------------------------------------
\subsection{Verificações adiadas para a fase do LLM}

\begin{itemize}
  \item \texttt{[3.x]} Compatibilidade tese × agente × período (marcos temporais)
  \item \texttt{[2.8]} Distinção entre as três situações de PPP não apresentado
  \item \texttt{[2.45]} Cabimento de decadência/prescrição com base em datas e tipo de ação
  \item \texttt{[1.4]} Cobertura dos períodos requeridos na petição inicial
\end{itemize}



\clearpage
\lhead{Contestação 12}
\section{Contestação 12 --- LUIZA HELENA BARROS}
% -----------------------------------------------------------------------
% Cabeçalho
% -----------------------------------------------------------------------
\begin{center}
  {\small Processo n.º 5013344-10.2025.4.03.6183 \quad NB: — \quad DER: 15/05/2024}
\end{center}

\vspace{0.5em}

\noindent\colorbox{corFundoReprovado}{%
  \begin{minipage}{\dimexpr\linewidth-2\fboxsep}
    \centering
    \vspace{0.3em}
    {\large \textcolor{corERRO}{\textbf{VEREDITO: REPROVADO}}}\\
    Erros: \textbf{1} \quad|\quad Pontos de verificação manual: \textbf{0}
    \vspace{0.3em}
  \end{minipage}%
}

\vspace{1em}


% -----------------------------------------------------------------------
% Bloco 1 - Estrutural
% -----------------------------------------------------------------------
\subsection{Bloco 1 --- Estrutural}

\statusINFO{Tipo de endereçamento (JEF x Justiça Federal comum)}\\
Endereçamento identificado: \textbf{JEF (Juizado Especial Federal)}.\\
\textit{Evidência: ... SENHOR(A) DOUTOR(A) JUIZ(A) FEDERAL DA 1ª VARA FEDERAL DO JUIZADO ESPECIAL FEDERAL DA SUBSEÇÃO JUDICIÁRIA DE SÃO PAULO - SP Processo nº 501334...}

\medskip
\statusERRO{Campo obrigatório: NB}\\
\textbf{NB ausente na minuta.}

\medskip
\statusOK{Campo obrigatório: DER} --- DER presente.

\medskip
\statusOK{Campo obrigatório: Número do processo} --- Número do processo presente.

\medskip
\statusOK{Campo obrigatório: Nome do autor} --- Nome do autor presente.\\
\textit{Evidência: ... a federal já qualificada nos autos do processo em epígrafe, que lhe move LUIZA HELENA BARROS, vem, respeitosamente, à presença de Vossa Excelência, por...}

\medskip
\statusOK{Campo obrigatório: Tipo de benefício} --- Tipo de benefício presente.

\medskip
\statusOK{Marcações de edição não removidas (colchetes, placeholders, instruções, amarelo)}\\
Nenhuma marcação de edição remanescente detectada.

% -----------------------------------------------------------------------
% Bloco 2 - Preliminares
% -----------------------------------------------------------------------
\subsection{Bloco 2 --- Preliminares (presença)}

\begin{longtable}{p{8.5cm} p{4.5cm}}
  \toprule
  \textbf{Preliminar} & \textbf{Status} \\
  \midrule
  \endfirsthead
  \toprule
  \textbf{Preliminar} & \textbf{Status} \\
  \midrule
  \endhead
  Juízo 100\% Digital                           & \textcolor{corOK}{\textbf{PRESENTE}} \\
  Audiência de Conciliação                      & \textcolor{corOK}{\textbf{PRESENTE}} \\
  Renúncia aos 60 Salários-Mínimos              & \textcolor{corOK}{\textbf{PRESENTE}} \\
  Decadência                                    & AUSENTE \\
  Prescrição                                    & AUSENTE \\
  Coisa Julgada                                 & AUSENTE \\
  Litispendência                                & AUSENTE \\
  PPP não apresentado administrativamente       & \textcolor{corOK}{\textbf{PRESENTE}} \\
  Períodos reconhecidos administrativamente     & AUSENTE \\
  Petição inicial inepta                        & AUSENTE \\
  \bottomrule
\end{longtable}

\statusOK{Renúncia aos 60 SM obrigatória quando o juízo é o JEF}\\
Processo no JEF e preliminar de Renúncia aos 60 SM presente.

% -----------------------------------------------------------------------
% Bloco 6 - Redistribuição
% -----------------------------------------------------------------------
\subsection{Bloco 6 --- Redistribuição}

\statusOK{Indício de categoria vigilante/vigia/guarda/policial} --- Sem indício detectado.

\medskip
\statusOK{Indício de profissional/ambiente de saúde + agente biológico} --- Sem indício detectado.

\medskip
\statusOK{Indício de professor com período posterior a 1981} --- Sem indício detectado.

% -----------------------------------------------------------------------
% Teses
% -----------------------------------------------------------------------
\subsection{Teses presentes (Blocos 3/4/5 --- catalogação)}

\statusINFO{Catalogação das teses padronizadas presentes (CTN-/DIVESP- etc.)}\\
Nenhuma tese padronizada (CTN-/DIVESP-) encontrada no texto.

% -----------------------------------------------------------------------
% Verificações adiadas
% -----------------------------------------------------------------------
\subsection{Verificações adiadas para a fase do LLM}

\begin{itemize}
  \item \texttt{[3.x]} Compatibilidade tese × agente × período (marcos temporais)
  \item \texttt{[2.8]} Distinção entre as três situações de PPP não apresentado
  \item \texttt{[2.45]} Cabimento de decadência/prescrição com base em datas e tipo de ação
  \item \texttt{[1.4]} Cobertura dos períodos requeridos na petição inicial
\end{itemize}



\clearpage
\lhead{Contestação 13}
\section{Contestação 13 --- MARCOS VINÍCIUS TEIXEIRA}
% -----------------------------------------------------------------------
% Cabeçalho
% -----------------------------------------------------------------------
\begin{center}
  {\small Processo n.º 5014455-21.2025.4.03.6301 \quad NB: 46/101.202.303-6 \quad DER: 10/07/2024}
\end{center}

\vspace{0.5em}

\noindent\colorbox{corFundoAprovado}{%
  \begin{minipage}{\dimexpr\linewidth-2\fboxsep}
    \centering
    \vspace{0.3em}
    {\large \textcolor{corOK}{\textbf{VEREDITO: APROVADO}}}\\
    Erros: \textbf{0} \quad|\quad Pontos de verificação manual: \textbf{0}
    \vspace{0.3em}
  \end{minipage}%
}

\vspace{1em}


% -----------------------------------------------------------------------
% Bloco 1 - Estrutural
% -----------------------------------------------------------------------
\subsection{Bloco 1 --- Estrutural}

\statusINFO{Tipo de endereçamento (JEF x Justiça Federal comum)}\\
Endereçamento identificado: \textbf{JEF (Juizado Especial Federal)}.\\
\textit{Evidência: ... SENHOR(A) DOUTOR(A) JUIZ(A) FEDERAL DA 8ª VARA FEDERAL DO JUIZADO ESPECIAL FEDERAL DA SUBSEÇÃO JUDICIÁRIA DE SÃO PAULO - SP Processo nº 501445...}

\medskip
\statusOK{Campo obrigatório: NB} --- NB presente.

\medskip
\statusOK{Campo obrigatório: DER} --- DER presente.

\medskip
\statusOK{Campo obrigatório: Número do processo} --- Número do processo presente.

\medskip
\statusOK{Campo obrigatório: Nome do autor} --- Nome do autor presente.\\
\textit{Evidência: ... a federal já qualificada nos autos do processo em epígrafe, que lhe move MARCOS VINÍCIUS TEIXEIRA, vem, respeitosamente, à presença de Vossa Excelência, por...}

\medskip
\statusOK{Campo obrigatório: Tipo de benefício} --- Tipo de benefício presente.

\medskip
\statusOK{Marcações de edição não removidas (colchetes, placeholders, instruções, amarelo)}\\
Nenhuma marcação de edição remanescente detectada.

% -----------------------------------------------------------------------
% Bloco 2 - Preliminares
% -----------------------------------------------------------------------
\subsection{Bloco 2 --- Preliminares (presença)}

\begin{longtable}{p{8.5cm} p{4.5cm}}
  \toprule
  \textbf{Preliminar} & \textbf{Status} \\
  \midrule
  \endfirsthead
  \toprule
  \textbf{Preliminar} & \textbf{Status} \\
  \midrule
  \endhead
  Juízo 100\% Digital                           & \textcolor{corOK}{\textbf{PRESENTE}} \\
  Audiência de Conciliação                      & \textcolor{corOK}{\textbf{PRESENTE}} \\
  Renúncia aos 60 Salários-Mínimos              & \textcolor{corOK}{\textbf{PRESENTE}} \\
  Decadência                                    & AUSENTE \\
  Prescrição                                    & AUSENTE \\
  Coisa Julgada                                 & AUSENTE \\
  Litispendência                                & AUSENTE \\
  PPP não apresentado administrativamente       & AUSENTE \\
  Períodos reconhecidos administrativamente     & AUSENTE \\
  Petição inicial inepta                        & AUSENTE \\
  \bottomrule
\end{longtable}

\statusOK{Renúncia aos 60 SM obrigatória quando o juízo é o JEF}\\
Processo no JEF e preliminar de Renúncia aos 60 SM presente.

% -----------------------------------------------------------------------
% Bloco 6 - Redistribuição
% -----------------------------------------------------------------------
\subsection{Bloco 6 --- Redistribuição}

\statusOK{Indício de categoria vigilante/vigia/guarda/policial} --- Sem indício detectado.

\medskip
\statusOK{Indício de profissional/ambiente de saúde + agente biológico} --- Sem indício detectado.

\medskip
\statusOK{Indício de professor com período posterior a 1981} --- Sem indício detectado.

% -----------------------------------------------------------------------
% Teses
% -----------------------------------------------------------------------
\subsection{Teses presentes (Blocos 3/4/5 --- catalogação)}

\statusINFO{Catalogação das teses padronizadas presentes (CTN-/DIVESP- etc.)}\\
Nenhuma tese padronizada (CTN-/DIVESP-) encontrada no texto.

% -----------------------------------------------------------------------
% Verificações adiadas
% -----------------------------------------------------------------------
\subsection{Verificações adiadas para a fase do LLM}

\begin{itemize}
  \item \texttt{[3.x]} Compatibilidade tese × agente × período (marcos temporais)
  \item \texttt{[2.8]} Distinção entre as três situações de PPP não apresentado
  \item \texttt{[2.45]} Cabimento de decadência/prescrição com base em datas e tipo de ação
  \item \texttt{[1.4]} Cobertura dos períodos requeridos na petição inicial
\end{itemize}



\clearpage
\lhead{Contestação 14}
\section{Contestação 14 --- RAIMUNDO NONATO SILVA}
% -----------------------------------------------------------------------
% Cabeçalho
% -----------------------------------------------------------------------
\begin{center}
  {\small Processo n.º 5015566-32.2025.4.03.6183 \quad NB: 46/121.232.343-9 \quad DER: 03/03/2024}
\end{center}

\vspace{0.5em}

\noindent\colorbox{corFundoAprovado}{%
  \begin{minipage}{\dimexpr\linewidth-2\fboxsep}
    \centering
    \vspace{0.3em}
    {\large \textcolor{corOK}{\textbf{VEREDITO: APROVADO}}}\\
    Erros: \textbf{0} \quad|\quad Pontos de verificação manual: \textbf{0}
    \vspace{0.3em}
  \end{minipage}%
}

\vspace{1em}


% -----------------------------------------------------------------------
% Bloco 1 - Estrutural
% -----------------------------------------------------------------------
\subsection{Bloco 1 --- Estrutural}

\statusINFO{Tipo de endereçamento (JEF x Justiça Federal comum)}\\
Endereçamento identificado: \textbf{JEF (Juizado Especial Federal)}.\\
\textit{Evidência: ... SENHOR(A) DOUTOR(A) JUIZ(A) FEDERAL DA 3ª VARA FEDERAL DO JUIZADO ESPECIAL FEDERAL DA SUBSEÇÃO JUDICIÁRIA DE SÃO PAULO - SP Processo nº 501556...}

\medskip
\statusOK{Campo obrigatório: NB} --- NB presente.

\medskip
\statusOK{Campo obrigatório: DER} --- DER presente.

\medskip
\statusOK{Campo obrigatório: Número do processo} --- Número do processo presente.

\medskip
\statusOK{Campo obrigatório: Nome do autor} --- Nome do autor presente.\\
\textit{Evidência: ... a federal já qualificada nos autos do processo em epígrafe, que lhe move RAIMUNDO NONATO SILVA, vem, respeitosamente, à presença de Vossa Excelência, por...}

\medskip
\statusOK{Campo obrigatório: Tipo de benefício} --- Tipo de benefício presente.

\medskip
\statusOK{Marcações de edição não removidas (colchetes, placeholders, instruções, amarelo)}\\
Nenhuma marcação de edição remanescente detectada.

% -----------------------------------------------------------------------
% Bloco 2 - Preliminares
% -----------------------------------------------------------------------
\subsection{Bloco 2 --- Preliminares (presença)}

\begin{longtable}{p{8.5cm} p{4.5cm}}
  \toprule
  \textbf{Preliminar} & \textbf{Status} \\
  \midrule
  \endfirsthead
  \toprule
  \textbf{Preliminar} & \textbf{Status} \\
  \midrule
  \endhead
  Juízo 100\% Digital                           & \textcolor{corOK}{\textbf{PRESENTE}} \\
  Audiência de Conciliação                      & \textcolor{corOK}{\textbf{PRESENTE}} \\
  Renúncia aos 60 Salários-Mínimos              & \textcolor{corOK}{\textbf{PRESENTE}} \\
  Decadência                                    & AUSENTE \\
  Prescrição                                    & AUSENTE \\
  Coisa Julgada                                 & AUSENTE \\
  Litispendência                                & AUSENTE \\
  PPP não apresentado administrativamente       & AUSENTE \\
  Períodos reconhecidos administrativamente     & AUSENTE \\
  Petição inicial inepta                        & AUSENTE \\
  \bottomrule
\end{longtable}

\statusOK{Renúncia aos 60 SM obrigatória quando o juízo é o JEF}\\
Processo no JEF e preliminar de Renúncia aos 60 SM presente.

% -----------------------------------------------------------------------
% Bloco 6 - Redistribuição
% -----------------------------------------------------------------------
\subsection{Bloco 6 --- Redistribuição}

\statusOK{Indício de categoria vigilante/vigia/guarda/policial} --- Sem indício detectado.

\medskip
\statusOK{Indício de profissional/ambiente de saúde + agente biológico} --- Sem indício detectado.

\medskip
\statusOK{Indício de professor com período posterior a 1981} --- Sem indício detectado.

% -----------------------------------------------------------------------
% Teses
% -----------------------------------------------------------------------
\subsection{Teses presentes (Blocos 3/4/5 --- catalogação)}

\statusINFO{Catalogação das teses padronizadas presentes (CTN-/DIVESP- etc.)}\\
Nenhuma tese padronizada (CTN-/DIVESP-) encontrada no texto.

% -----------------------------------------------------------------------
% Verificações adiadas
% -----------------------------------------------------------------------
\subsection{Verificações adiadas para a fase do LLM}

\begin{itemize}
  \item \texttt{[3.x]} Compatibilidade tese × agente × período (marcos temporais)
  \item \texttt{[2.8]} Distinção entre as três situações de PPP não apresentado
  \item \texttt{[2.45]} Cabimento de decadência/prescrição com base em datas e tipo de ação
  \item \texttt{[1.4]} Cobertura dos períodos requeridos na petição inicial
\end{itemize}



\clearpage
\lhead{Contestação 15}
\section{Contestação 15 --- TEREZA CRISTINA LOPES}
% -----------------------------------------------------------------------
% Cabeçalho
% -----------------------------------------------------------------------
\begin{center}
  {\small Processo n.º 5016677-43.2025.4.03.6301 \quad NB: 46/131.242.353-4 \quad DER: 25/04/2024}
\end{center}

\vspace{0.5em}

\noindent\colorbox{corFundoReprovado}{%
  \begin{minipage}{\dimexpr\linewidth-2\fboxsep}
    \centering
    \vspace{0.3em}
    {\large \textcolor{corERRO}{\textbf{VEREDITO: REPROVADO}}}\\
    Erros: \textbf{1} \quad|\quad Pontos de verificação manual: \textbf{0}
    \vspace{0.3em}
  \end{minipage}%
}

\vspace{1em}


% -----------------------------------------------------------------------
% Bloco 1 - Estrutural
% -----------------------------------------------------------------------
\subsection{Bloco 1 --- Estrutural}

\statusINFO{Tipo de endereçamento (JEF x Justiça Federal comum)}\\
Endereçamento identificado: \textbf{JEF (Juizado Especial Federal)}.\\
\textit{Evidência: ... SENHOR(A) DOUTOR(A) JUIZ(A) FEDERAL DA 6ª VARA FEDERAL DO JUIZADO ESPECIAL FEDERAL DA SUBSEÇÃO JUDICIÁRIA DE SÃO PAULO - SP Processo nº 501667...}

\medskip
\statusOK{Campo obrigatório: NB} --- NB presente.

\medskip
\statusOK{Campo obrigatório: DER} --- DER presente.

\medskip
\statusOK{Campo obrigatório: Número do processo} --- Número do processo presente.

\medskip
\statusOK{Campo obrigatório: Nome do autor} --- Nome do autor presente.\\
\textit{Evidência: ... a federal já qualificada nos autos do processo em epígrafe, que lhe move TEREZA CRISTINA LOPES, vem, respeitosamente, à presença de Vossa Excelência, por...}

\medskip
\statusOK{Campo obrigatório: Tipo de benefício} --- Tipo de benefício presente.

\medskip
\statusOK{Marcações de edição não removidas (colchetes, placeholders, instruções, amarelo)}\\
Nenhuma marcação de edição remanescente detectada.

% -----------------------------------------------------------------------
% Bloco 2 - Preliminares
% -----------------------------------------------------------------------
\subsection{Bloco 2 --- Preliminares (presença)}

\begin{longtable}{p{8.5cm} p{4.5cm}}
  \toprule
  \textbf{Preliminar} & \textbf{Status} \\
  \midrule
  \endfirsthead
  \toprule
  \textbf{Preliminar} & \textbf{Status} \\
  \midrule
  \endhead
  Juízo 100\% Digital                           & AUSENTE \\
  Audiência de Conciliação                      & \textcolor{corOK}{\textbf{PRESENTE}} \\
  Renúncia aos 60 Salários-Mínimos              & AUSENTE \\
  Decadência                                    & AUSENTE \\
  Prescrição                                    & AUSENTE \\
  Coisa Julgada                                 & AUSENTE \\
  Litispendência                                & AUSENTE \\
  PPP não apresentado administrativamente       & AUSENTE \\
  Períodos reconhecidos administrativamente     & AUSENTE \\
  Petição inicial inepta                        & AUSENTE \\
  \bottomrule
\end{longtable}

\statusERRO{Renúncia aos 60 SM obrigatória quando o juízo é o JEF}\\
Processo no JEF, mas a preliminar de Renúncia aos 60 salários-mínimos está AUSENTE (art. 3º, §2º, Lei 10.259/2001).

% -----------------------------------------------------------------------
% Bloco 6 - Redistribuição
% -----------------------------------------------------------------------
\subsection{Bloco 6 --- Redistribuição}

\statusOK{Indício de categoria vigilante/vigia/guarda/policial} --- Sem indício detectado.

\medskip
\statusOK{Indício de profissional/ambiente de saúde + agente biológico} --- Sem indício detectado.

\medskip
\statusOK{Indício de professor com período posterior a 1981} --- Sem indício detectado.

% -----------------------------------------------------------------------
% Teses
% -----------------------------------------------------------------------
\subsection{Teses presentes (Blocos 3/4/5 --- catalogação)}

\statusINFO{Catalogação das teses padronizadas presentes (CTN-/DIVESP- etc.)}\\
Nenhuma tese padronizada (CTN-/DIVESP-) encontrada no texto.

% -----------------------------------------------------------------------
% Verificações adiadas
% -----------------------------------------------------------------------
\subsection{Verificações adiadas para a fase do LLM}

\begin{itemize}
  \item \texttt{[3.x]} Compatibilidade tese × agente × período (marcos temporais)
  \item \texttt{[2.8]} Distinção entre as três situações de PPP não apresentado
  \item \texttt{[2.45]} Cabimento de decadência/prescrição com base em datas e tipo de ação
  \item \texttt{[1.4]} Cobertura dos períodos requeridos na petição inicial
\end{itemize}



\clearpage
\lhead{Contestação 16}
\section{Contestação 16 --- EDSON LUÍS CAVALCANTE}
% -----------------------------------------------------------------------
% Cabeçalho
% -----------------------------------------------------------------------
\begin{center}
  {\small Processo n.º 5017788-54.2025.4.03.6183 \quad NB: 46/141.252.363-1 \quad DER: 30/09/2023}
\end{center}

\vspace{0.5em}

\noindent\colorbox{corFundoAprovado}{%
  \begin{minipage}{\dimexpr\linewidth-2\fboxsep}
    \centering
    \vspace{0.3em}
    {\large \textcolor{corOK}{\textbf{VEREDITO: APROVADO}}}\\
    Erros: \textbf{0} \quad|\quad Pontos de verificação manual: \textbf{0}
    \vspace{0.3em}
  \end{minipage}%
}

\vspace{1em}


% -----------------------------------------------------------------------
% Bloco 1 - Estrutural
% -----------------------------------------------------------------------
\subsection{Bloco 1 --- Estrutural}

\statusINFO{Tipo de endereçamento (JEF x Justiça Federal comum)}\\
Endereçamento identificado: \textbf{JEF (Juizado Especial Federal)}.\\
\textit{Evidência: ... SENHOR(A) DOUTOR(A) JUIZ(A) FEDERAL DA 10ª VARA FEDERAL DO JUIZADO ESPECIAL FEDERAL DA SUBSEÇÃO JUDICIÁRIA DE SÃO PAULO - SP Processo nº 501778...}

\medskip
\statusOK{Campo obrigatório: NB} --- NB presente.

\medskip
\statusOK{Campo obrigatório: DER} --- DER presente.

\medskip
\statusOK{Campo obrigatório: Número do processo} --- Número do processo presente.

\medskip
\statusOK{Campo obrigatório: Nome do autor} --- Nome do autor presente.\\
\textit{Evidência: ... a federal já qualificada nos autos do processo em epígrafe, que lhe move EDSON LUÍS CAVALCANTE, vem, respeitosamente, à presença de Vossa Excelência, por...}

\medskip
\statusOK{Campo obrigatório: Tipo de benefício} --- Tipo de benefício presente.

\medskip
\statusOK{Marcações de edição não removidas (colchetes, placeholders, instruções, amarelo)}\\
Nenhuma marcação de edição remanescente detectada.

% -----------------------------------------------------------------------
% Bloco 2 - Preliminares
% -----------------------------------------------------------------------
\subsection{Bloco 2 --- Preliminares (presença)}

\begin{longtable}{p{8.5cm} p{4.5cm}}
  \toprule
  \textbf{Preliminar} & \textbf{Status} \\
  \midrule
  \endfirsthead
  \toprule
  \textbf{Preliminar} & \textbf{Status} \\
  \midrule
  \endhead
  Juízo 100\% Digital                           & \textcolor{corOK}{\textbf{PRESENTE}} \\
  Audiência de Conciliação                      & \textcolor{corOK}{\textbf{PRESENTE}} \\
  Renúncia aos 60 Salários-Mínimos              & \textcolor{corOK}{\textbf{PRESENTE}} \\
  Decadência                                    & AUSENTE \\
  Prescrição                                    & AUSENTE \\
  Coisa Julgada                                 & AUSENTE \\
  Litispendência                                & AUSENTE \\
  PPP não apresentado administrativamente       & AUSENTE \\
  Períodos reconhecidos administrativamente     & AUSENTE \\
  Petição inicial inepta                        & AUSENTE \\
  \bottomrule
\end{longtable}

\statusOK{Renúncia aos 60 SM obrigatória quando o juízo é o JEF}\\
Processo no JEF e preliminar de Renúncia aos 60 SM presente.

% -----------------------------------------------------------------------
% Bloco 6 - Redistribuição
% -----------------------------------------------------------------------
\subsection{Bloco 6 --- Redistribuição}

\statusOK{Indício de categoria vigilante/vigia/guarda/policial} --- Sem indício detectado.

\medskip
\statusOK{Indício de profissional/ambiente de saúde + agente biológico} --- Sem indício detectado.

\medskip
\statusOK{Indício de professor com período posterior a 1981} --- Sem indício detectado.

% -----------------------------------------------------------------------
% Teses
% -----------------------------------------------------------------------
\subsection{Teses presentes (Blocos 3/4/5 --- catalogação)}

\statusINFO{Catalogação das teses padronizadas presentes (CTN-/DIVESP- etc.)}\\
Nenhuma tese padronizada (CTN-/DIVESP-) encontrada no texto.

% -----------------------------------------------------------------------
% Verificações adiadas
% -----------------------------------------------------------------------
\subsection{Verificações adiadas para a fase do LLM}

\begin{itemize}
  \item \texttt{[3.x]} Compatibilidade tese × agente × período (marcos temporais)
  \item \texttt{[2.8]} Distinção entre as três situações de PPP não apresentado
  \item \texttt{[2.45]} Cabimento de decadência/prescrição com base em datas e tipo de ação
  \item \texttt{[1.4]} Cobertura dos períodos requeridos na petição inicial
\end{itemize}



\clearpage
\lhead{Gabarito e divergências}
\section{Gabarito de referência e análise das divergências}

Esta seção confronta o veredito emitido pela ferramenta (\textbf{Fase~1 ---
camada determinística, sem LLM}) com o \textbf{gabarito} de referência das 14
contestações e explica, caso a caso, \emph{por que} a ferramenta diverge.

\medskip
\noindent\textbf{Ponto central.} A Fase~1 só executa verificações de
\emph{presença/ausência} e de \emph{padrão textual}. Ela \textbf{não} avalia
compatibilidade entre tese, agente nocivo e período (marcos temporais), nem
coteja a peça com os pedidos da petição inicial ou com o formulário de
contexto/PA --- essas checagens estão declaradas, em todos os relatórios, como
\textbf{``Verificações adiadas para a fase do LLM''} (\texttt{[3.x]},
\texttt{[2.8]}, \texttt{[2.45]}, \texttt{[1.4]}). Portanto, as divergências
abaixo \textbf{não são falhas da camada determinística}: são exatamente os casos
que ela, por projeto, deixa para a camada de raciocínio. Vale registrar que
\textbf{nenhum} \status{ERRO} emitido é indevido --- toda reprovação da
ferramenta (05, 10, 12, 15) coincide com o gabarito (precisão de 100\% nos erros
apontados).

\subsection*{Comparativo: ferramenta (Fase 1) $\times$ gabarito}

\begin{longtable}{c c c c p{7.3cm}}
  \toprule
  \textbf{Nº} & \textbf{Fase 1} & \textbf{Gabarito} & \textbf{Converge?} & \textbf{Observação / motivo da divergência} \\
  \midrule
  \endfirsthead
  \toprule
  \textbf{Nº} & \textbf{Fase 1} & \textbf{Gabarito} & \textbf{Converge?} & \textbf{Observação / motivo da divergência} \\
  \midrule
  \endhead
  03 & \textcolor{corOK}{APROVADO}  & APROVAR  & \textcolor{corOK}{Sim} & --- \\
  04 & \textcolor{corOK}{APROVADO}  & APROVAR  & \textcolor{corOK}{Sim} & --- \\
  05 & \textcolor{corERRO}{REPROVADO} & REPROVAR & \textcolor{corOK}{Sim} & Anotação de edição não removida (detecção determinística). \\
  06 & \textcolor{corOK}{APROVADO}  & REPROVAR & \textcolor{corERRO}{\textbf{Não}} & Limiar de ruído (90 vs \textbf{85 dB(A) NEN}), metodologia (NR-15 vs \textbf{NEN/NHO-01}) e vedação pós-13/11/2019 ausente --- exigem análise agente$\times$período (\texttt{[3.x]}, \texttt{[1.4]}). \\
  07 & \textcolor{corOK}{APROVADO}  & APROVAR  & \textcolor{corOK}{Sim} & --- \\
  08 & \textcolor{corOK}{APROVADO}  & REPROVAR & \textcolor{corVERIFICAR}{\textbf{Parcial}} & \textbf{VIGILANTE} sinalizado (flag de redistribuição), mas a regra do Bloco~6 é \emph{não-bloqueante} por projeto. \\
  09 & \textcolor{corOK}{APROVADO}  & REPROVAR & \textcolor{corVERIFICAR}{\textbf{Parcial}} & \textbf{SAÚDE}/agente biológico sinalizado (flag), também não-bloqueante por projeto. \\
  10 & \textcolor{corERRO}{REPROVADO} & REPROVAR & \textcolor{corOK}{Sim*} & Reprovada por \textbf{Tipo de benefício ausente} (determinístico); \textbf{PROFESSOR} pós-1981 também foi sinalizado. \\
  11 & \textcolor{corOK}{APROVADO}  & REPROVAR & \textcolor{corERRO}{\textbf{Não}} & Indeferimento forçado (tempo especial = Não + procuração \emph{ad judicia}) --- depende do PA/contexto, fora do texto da minuta. \\
  12 & \textcolor{corERRO}{REPROVADO} & REPROVAR & \textcolor{corOK}{Sim*} & Reprovada por \textbf{NB ausente} (determinístico); ``sem pedido de tempo especial'' é da camada LLM. \\
  13 & \textcolor{corOK}{APROVADO}  & REPROVAR & \textcolor{corERRO}{\textbf{Não}} & Impugna período não requerido e omite período requerido --- exige cotejo com a inicial (\texttt{[1.4]}). \\
  14 & \textcolor{corOK}{APROVADO}  & REPROVAR & \textcolor{corERRO}{\textbf{Não}} & Vício CREA/CRM obsoleto + limiar de calor (24,7 vs \textbf{25 ºC}) --- exigem base de teses e marcos temporais (\texttt{[3.x]}). \\
  15 & \textcolor{corERRO}{REPROVADO} & REPROVAR & \textcolor{corOK}{Sim} & Faltam \textbf{Juízo 100\% Digital} e \textbf{Renúncia 60SM} no JEF (determinístico). \\
  16 & \textcolor{corOK}{APROVADO}  & APROVAR  & \textcolor{corOK}{Sim} & --- \\
  \bottomrule
\end{longtable}

\noindent\footnotesize{* ``Sim*'': o veredito coincide com o gabarito, porém a
reprovação da ferramenta decorre de um erro \emph{estrutural} detectável
(NB/benefício), enquanto o motivo apontado no gabarito depende, em parte, da
camada LLM.}\normalsize

\subsection*{Por que a Fase 1 não reprova esses casos (por grupo)}

\begin{itemize}[leftmargin=1.4em]
  \item \textbf{Grupo A --- compatibilidade tese $\times$ agente $\times$ período
        (\texttt{[3.x]}, \texttt{[1.4]}): casos 06 e 14.} Detectar que ``90 dB(A)''
        está errado depois de 19/11/2003 (correto: 85 dB(A) NEN), que a
        metodologia deveria ser NEN/NHO-01, que o limiar de calor muda de 25 ºC
        para 24,7 ºC após a Portaria SEPRT 1.359/2019, ou que o vício CREA/CRM
        está obsoleto, exige conhecer marcos temporais e uma base de teses
        válidas. A camada determinística apenas cataloga teses; não julga
        adequação.
  \item \textbf{Grupo B --- cotejo com a petição inicial (\texttt{[1.4]}): caso 13
        (e a vedação ausente do 06).} Saber que a peça impugnou um período que o
        autor não pediu --- ou omitiu um que ele pediu --- requer comparar a
        contestação com os períodos da inicial, o que a Fase~1 não faz.
  \item \textbf{Grupo C --- contexto/PA fora do texto da minuta: caso 11.} O
        gatilho de \emph{indeferimento forçado} (``possui tempo especial? Não'' +
        procuração \emph{ad judicia} + ausência de análise técnica) vem do
        formulário do processo administrativo, não de um padrão textual da
        minuta.
  \item \textbf{Grupo D --- redistribuição sinalizada, mas não-bloqueante: casos
        08 e 09 (e reforço em 10).} O Bloco~6 identifica corretamente os perfis
        VIGILANTE, SAÚDE e PROFESSOR e emite a \emph{flag} de redistribuição,
        porém --- por decisão de projeto --- essa sinalização é um alerta
        (\status{INFO}/\status{VERIFICAR}), nunca um \status{ERRO}. Logo, o
        veredito permanece \textsc{aprovado} ainda que o caso deva ser
        redistribuído.
\end{itemize}

\subsection*{Síntese}

Das 14 contestações, o veredito da Fase~1 \textbf{coincide com o gabarito em 8}
(03, 04, 05, 07, 10, 12, 15, 16); em \textbf{2} a ferramenta \emph{sinaliza} o
problema mas não o converte em reprovação por design (08, 09); e em \textbf{4} a
detecção foi legitimamente adiada para a camada de LLM (06, 11, 13, 14).
Nenhuma reprovação foi indevida. As divergências, portanto, medem exatamente a
fronteira entre a camada determinística já entregue e a camada de raciocínio
ainda pendente.

\end{document}
